%% 28장 큰 고유값 문제
\chapter{큰 고유값 문제}
\label{ch:large-eigen}

\abstract{23장의 QR 알고리즘은 행렬 전체를 Hessenberg 꼴로 바꾸므로 $O(n^3)$의 일과 $O(n^2)$의 저장이 든다. $n$이 수십만인 희소 행렬에서 몇 개의 고유값만 필요하다면 그 비용은 감당할 수 없다. 이 장에서는 앞 장의 Krylov 부분공간을 고유값 문제에 쓴다. Krylov 부분공간에 행렬을 사영한 작은 행렬의 고유값, 곧 Ritz 값이 스펙트럼의 바깥쪽 고유값으로 빠르게 수렴하는 이유를 Chebyshev 다항식으로 설명하고, 유한 정밀도의 Lanczos 과정이 같은 고유값을 여러 번 찾아내는 현상을 직접 계산으로 본다. 기저가 커지는 문제는 원하지 않는 Ritz 값을 이동으로 쓰는 암시적 재시작으로 해결하는데, 그 정당성은 23장의 암시적 Q 정리에서 나온다. 스펙트럼 안쪽의 고유값을 위한 스펙트럼 변환, 특이값을 위한 Golub--Kahan--Lanczos 이중대각화를 다루고, stardust가 Accelerate에 없는 ARPACK을 서브루틴 단위로 옮긴 방식을 본다.}

\section{Krylov 부분공간의 Rayleigh--Ritz}
\label{sec:ritz}

\subsection{Ritz 값과 잔차}

정규직교 열의 행렬 $V_k\in\mathbb R^{n\times k}$가 주어지면, $A$의 고유값을 $V_k$의 열공간 안에서 근사하는 자연스러운 방법은 사영 $H_k=V_k^{\mathsf T}AV_k$의 고유쌍 $H_ks=\theta s$를 구해 $(\theta,V_ks)$를 근사 고유쌍으로 쓰는 것이다. 이것을 \term{Rayleigh--Ritz 방법}{Rayleigh--Ritz method}이라 하고, $\theta$를 \term{Ritz 값}{Ritz value}, $y=V_ks$를 \term{Ritz 벡터}{Ritz vector}라 한다. 잔차 $Ay-\theta y$는 $V_k$의 열공간에 직교하므로 이 방법은 앞 장의 Galerkin 조건을 고유값 문제에 쓴 것이다.

$V_k$가 Krylov 부분공간의 Arnoldi 기저이면 Ritz 쌍의 잔차를 공짜로 얻는다. Arnoldi 관계 $AV_k=V_kH_k+h_{k+1,k}v_{k+1}e_k^{\mathsf T}$에 $s$를 곱하면
\begin{equation}
\label{eq:ritz-residual}
Ay-\theta y=V_k(H_ks-\theta s)+h_{k+1,k}(e_k^{\mathsf T}s)v_{k+1}=h_{k+1,k}(e_k^{\mathsf T}s)v_{k+1}
\end{equation}
이고, 따라서 $\|Ay-\theta y\|=|h_{k+1,k}|\,|e_k^{\mathsf T}s|$이다. $n$차원의 벡터를 만들지 않고, 작은 고유벡터 $s$의 마지막 성분만으로 잔차를 안다. 대칭 행렬이면 이 잔차가 오차의 확실한 상계가 된다.

\begin{proposition}
\label{prop:ritz-residual}
$A$가 대칭이고 $\|y\|=1$이면, 어떤 고유값 $\lambda$가 $|\lambda-\theta|\le\|Ay-\theta y\|$를 만족한다.
\end{proposition}

\begin{proof}
$\theta$가 고유값이면 자명하다. 아니면 $A=Q\Lambda Q^{\mathsf T}$로 쓸 때 $y=(A-\theta I)^{-1}(Ay-\theta y)$이므로 $1=\|y\|\le\|(A-\theta I)^{-1}\|\,\|Ay-\theta y\|=\|Ay-\theta y\|/\min_i|\lambda_i-\theta|$이다.
\end{proof}

비대칭 행렬에서는 4장의 Bauer--Fike 정리처럼 고유벡터 행렬의 조건수가 곱해지므로, 작은 잔차가 반드시 정확한 고유값을 뜻하지는 않는다.

\subsection{바깥쪽 고유값의 수렴}

대칭 행렬에서 Lanczos 과정의 Ritz 값은 스펙트럼의 양 끝에서 먼저 수렴한다. 이유는 가장 큰 Ritz 값이 Krylov 부분공간 위의 Rayleigh 몫의 최댓값이라는 데 있다. 앞 장의 CG처럼 Krylov 부분공간의 원소는 $p(A)v$로 쓸 수 있으므로 이 최댓값의 크기는 다시 다항식 근사의 문제가 된다. 이 결과는 Kaniel이 처음 보이고 Paige와 Saad가 다듬었다\cite{kaniel1966,saad1980rates}.

\begin{theorem}[Kaniel--Paige--Saad]
\label{thm:kps}
대칭 행렬 $A$의 고유값을 $\lambda_1>\lambda_2\ge\dots\ge\lambda_n$, $\lambda_1$의 단위 고유벡터를 $q_1$이라 하자. $\mathcal K_k(A,v)$에서 얻은 가장 큰 Ritz 값 $\theta_1^{(k)}$는
\[
0\le\lambda_1-\theta_1^{(k)}\le(\lambda_1-\lambda_n)\Bigl(\frac{\tan\varphi_1}{T_{k-1}(1+2\gamma_1)}\Bigr)^2,\qquad\gamma_1=\frac{\lambda_1-\lambda_2}{\lambda_2-\lambda_n}
\]
를 만족한다. 여기서 $\varphi_1$은 $v$와 $q_1$ 사이의 각이고 $T_{k-1}$은 Chebyshev 다항식이다.
\end{theorem}

\begin{proof}
$\theta_1^{(k)}$는 $V_k^{\mathsf T}AV_k$의 가장 큰 고유값이므로, 4장의 Courant--Fischer 정리에 의해 $\mathcal K_k$ 위의 Rayleigh 몫 $R(x)=x^{\mathsf T}Ax/x^{\mathsf T}x$의 최댓값이다. 그래서 왼쪽 부등식이 성립하고, 오른쪽을 보이려면 $\mathcal K_k$의 원소 $x$ 하나에 대해 $\lambda_1-R(x)$를 위에서 누르면 된다. $v=\sum_ic_iq_i$로 쓰고 차수 $k-1$의 다항식 $p$에 대해 $x=p(A)v$로 두면
\[
\lambda_1-R(x)=\frac{\sum_{i\ge2}(\lambda_1-\lambda_i)p(\lambda_i)^2c_i^2}{\sum_ip(\lambda_i)^2c_i^2}\le(\lambda_1-\lambda_n)\max_{i\ge2}\frac{p(\lambda_i)^2}{p(\lambda_1)^2}\cdot\frac{\sum_{i\ge2}c_i^2}{c_1^2}
\]
이다. 분모에서 $i\ge2$의 항을 버렸다. 마지막 몫은 $\tan^2\varphi_1$이다. 이제 $[\lambda_n,\lambda_2]$를 $[-1,1]$로 보내는 $\ell(t)=1+2(t-\lambda_2)/(\lambda_2-\lambda_n)$을 두고 $p=T_{k-1}\circ\ell$로 고르면, $i\ge2$에서 $|p(\lambda_i)|\le1$이고 $p(\lambda_1)=T_{k-1}(1+2\gamma_1)$이다.
\end{proof}

앞 장의 \eqref{eq:chebyshev-outside}에 의해 $T_{k-1}(1+2\gamma_1)$은 대략 $\bigl(1+2\sqrt{\gamma_1}\bigr)^{k-1}$처럼 자라므로 오차는 걸음마다 $(1+2\sqrt{\gamma_1})^{-2}$배쯤 줄어든다. 23장의 거듭제곱법은 이동을 가장 잘 골라도 고유값의 오차를 걸음마다 $(1+2\gamma_1)^{-2}$배쯤 줄이는 것과 견주면, Lanczos는 상대 간격 $\gamma_1$ 대신 그 제곱근을 치른다. CG가 $\kappa$ 대신 $\sqrt\kappa$를 치른 것과 같은 이득이다.

이 상계도 CG의 상계처럼 비관적일 수 있다. $n=1000$이고 고유값이 $10$, $9.8$, 그리고 $[0,9]$의 균등 난수 $998$개인 대각 행렬에 무작위 시작 벡터로 Lanczos를 돌려 보자(직교성을 지키려고 매번 다시 직교화했다). $\gamma_1\approx0.0205$, $\tan\varphi_1\approx35.3$이다. $10$걸음에서 실제 오차 $\lambda_1-\theta_1$은 $6.2\times10^{-2}$이고 상계는 $2.9\times10^2$이며, $20$걸음에서는 $1.1\times10^{-5}$ 대 $0.98$, $30$걸음에서는 $2.1\times10^{-11}$ 대 $3.3\times10^{-3}$이다. 상계가 느린 것은 $\lambda_2=9.8$이 $\lambda_1$에 가까워 $\gamma_1$이 작기 때문인데, 실제 Lanczos는 $\lambda_2$도 함께 찾아내고 나면 그다음 고유값 $9$ 근처까지의 넓은 간격으로 수렴한다. 앞 장에서 본 초선형 수렴이 여기서도 일어난다.

반면 스펙트럼 안쪽의 고유값이나 바깥쪽이라도 서로 붙어 있는 고유값은 느리게 수렴한다. 1차원 Laplace 행렬 $\operatorname{tridiag}(-1,2,-1)$은 고유값이 $2-2\cos(j\pi/(n+1))$이어서, 가장 작은 고유값의 상대 간격이 $\gamma_1\approx3\pi^2/\bigl(4(n+1)^2\bigr)$ 정도로 작다. 이런 고유값을 어떻게 다루는지는 \cref{sec:spectral-transform}에서 본다.

\section{유한 정밀도의 Lanczos 과정}
\label{sec:lanczos-finite}

\subsection{가짜 고유값}

Lanczos 과정은 걸음마다 일이 $O(n)$이고 벡터 몇 개만 저장하면 되므로 원리적으로는 큰 대칭 행렬의 고유값을 구하는 이상적인 방법이다. Lanczos가 1950년에 이 방법을 발표했을 때도 그렇게 기대되었다\cite{lanczos1950}. 그러나 앞 장에서 보았듯이 유한 정밀도에서는 Lanczos 벡터의 직교성이 무너진다. 고유값 문제에서는 그 결과가 더 눈에 띈다. 이미 찾은 고유값이 Ritz 값으로 다시, 또다시 나타난다.

앞 장의 Strakoš 꼴 대각 행렬을 $n=100$, $\lambda_1=0.1$, $\lambda_n=100$, $\rho=0.9$로 만들고 $v=(1,\dots,1)^{\mathsf T}$에서 Lanczos를 돌려 보자. 다시 직교화하지 않으면 가장 큰 고유값 $100$은 $13$걸음에 상대 $10^{-8}$까지 수렴하고, 그 직후인 $21$걸음에 Lanczos 벡터들의 직교성이 $\max|V_k^{\mathsf T}V_k-I|>10^{-8}$로 무너지기 시작한다. 정확한 산술이라면 $100$걸음에 모든 고유값이 나와야 하지만, $100$걸음에서 $T_{100}$의 고유값 가운데 가장 큰 고유값의 복사본이 $5$개, 두 번째 고유값의 복사본도 $5$개이고, 상대 $10^{-8}$ 안에서 맞춘 고유값은 $100$개 가운데 $39$개뿐이다. $300$걸음까지 가면 가장 큰 고유값의 복사본은 $17$개로 늘고 찾은 고유값은 $62$개가 된다. 매번 완전히 다시 직교화하면 $100$걸음에 $100$개를 모두 찾는다.

이 복사본을 \term{가짜 고유값}{ghost eigenvalue}이라 부르지만 사실은 가짜가 아니라 진짜 고유값의 정확한 근사이다. 다만 그 중복도가 틀렸다. Paige는 1970년대의 일련의 연구로 이 현상을 해명했다\cite{paige1976error}. 직교성은 아무렇게나 무너지지 않고, Ritz 값 하나가 고유값에 수렴하면 바로 그 Ritz 벡터 방향으로 무너진다. 수렴한 고유벡터 방향의 성분이 반올림으로 다시 생겨나 자라면 Lanczos 과정은 그 방향을 처음 보는 것처럼 다시 찾는다. 위의 실험에서 직교성이 무너진 시점이 가장 큰 고유값이 수렴한 직후였던 것이 그 증거이다.

\subsection{다시 직교화하는 방법}

해결책은 직교성을 다시 세우는 것이다. 가장 확실한 방법은 새 벡터를 앞의 모든 Lanczos 벡터에 대해 다시 직교화하는 \term{완전 재직교화}{full reorthogonalization}인데, 그러면 걸음마다 일이 $O(nk)$로 늘고 모든 벡터를 저장해야 해서 Arnoldi 과정과 비용이 같아진다. Paige의 해명은 더 싼 길을 열었다. 직교성이 수렴한 Ritz 벡터 방향으로만 무너진다면 그 방향들에 대해서만 직교화하면 되고, 이 방법이 Parlett와 Scott의 \term{선택적 재직교화}{selective orthogonalization}이다\cite{parlett1979selective}. 반대로 재직교화를 아예 하지 않고 복사본을 사후에 가려내는 Cullum과 Willoughby의 방법도 있다\cite{cullum1985}.

재직교화를 한 번 하면 충분한가? 앞의 기저 벡터들에 대해 고전 Gram--Schmidt로 직교화한 벡터 $w$의 노름이 직교화 전보다 크게 줄었다면, 그것은 $w$가 거의 기저의 열공간 안에 있었다는 뜻이고 남은 작은 성분에는 상대적으로 큰 반올림이 섞여 있다. Daniel, Gragg, Kaufman, Stewart는 이때 한 번 더 직교화하는 기준을 세웠다\cite{daniel1976dgks}. 직교화 뒤의 노름이 직교화 전의 $\eta$배보다 작아졌을 때만 다시 하는데, 흔히 $\eta=1/\sqrt2\approx0.707$ 근처의 값을 쓰는데, ARPACK은 $0.717$을 쓴다. Kahan이 관찰하고 Parlett가 증명한 대로, 이렇게 두 번 하면 반올림 수준의 직교성을 얻기에 충분하다\cite{parlett1998}. 흔히 \term{두 번이면 충분하다}{twice is enough}고 줄여 말한다.

\section{재시작}
\label{sec:restart}

\subsection{다항식 거르개}

완전 재직교화를 하면 Lanczos도 Arnoldi처럼 $k$걸음에 $O(nk)$의 메모리를 쓰므로, 수렴이 느린 고유값을 기다리며 기저를 무한히 늘릴 수 없다. 앞 장의 재시작 GMRES처럼 $m$걸음마다 새로 시작해야 하는데, 문제는 무엇에서 다시 시작하느냐이다. 원하는 고유값이 $k$개라면 이상적인 시작 벡터는 그 $k$개의 고유벡터가 펼치는 공간 안에 있는 벡터이다. 그런 벡터에서 시작하면 Krylov 부분공간이 그 불변 부분공간에서 정확히 멈추기 때문이다. 그래서 시작 벡터 $v_1$을 $\psi(A)v_1$로 바꾸되, 다항식 $\psi$가 원하지 않는 고유값 근처에서 작아지도록 고른다. 이것을 \term{다항식 거르개}{polynomial filter}라 한다.

원하지 않는 고유값의 위치는 모르지만 그 근사는 지금의 Krylov 부분공간이 이미 알려 준다. $m$개의 Ritz 값 가운데 원하지 않는 $p=m-k$개 $\mu_1,\dots,\mu_p$를 근으로 하는 $\psi(t)=\prod_j(t-\mu_j)$를 쓰면 된다. 이렇게 고른 이동을 \term{정확한 이동}{exact shifts}이라 한다. 다만 $\psi(A)v_1$을 계산하려면 곱이 $p$번 더 들고 거기서 Arnoldi를 처음부터 다시 해야 한다.

\subsection{암시적 재시작}

Sorensen은 이 일을 곱 없이, 작은 행렬의 QR 걸음만으로 할 수 있음을 보였다\cite{sorensen1992implicit}. $m$걸음의 Arnoldi 관계
\[
AV_m=V_mH_m+f_me_m^{\mathsf T},\qquad f_m=h_{m+1,m}v_{m+1}
\]
에서 출발해 23장의 이동 QR 걸음을 이동 $\mu_1,\dots,\mu_p$로 $H_m$에 $p$번 적용한다.

\begin{theorem}[암시적 재시작]
\label{thm:implicit-restart}
$H_m-\mu_jI=Q_jR_j$, $H_m\gets Q_j^{\mathsf T}H_mQ_j$를 $j=1,\dots,p$에 대해 차례로 하고 $Q=Q_1\cdots Q_p$, $V^+=V_mQ$, $H^+=Q^{\mathsf T}H_mQ$로 두자. 그러면 $V^+$와 $H^+$의 첫 $k=m-p$열은
\[
AV_k^+=V_k^+H_k^++f_k^+e_k^{\mathsf T}
\]
꼴의 $k$걸음 Arnoldi 관계를 만족하고, 그 시작 벡터 $v_1^+$는 $\psi(A)v_1$의 방향이다. 따라서 이것은 $\psi(A)v_1$에서 시작해 $k$걸음을 한 Arnoldi 관계와 같다.
\end{theorem}

\begin{proof}
이동 하나에 대해 보자. $(A-\mu I)V_m=V_m(H_m-\mu I)+f_me_m^{\mathsf T}=V_mQ_1R_1+f_me_m^{\mathsf T}$이다. 첫 열을 보면, $m>1$이므로 $e_m^{\mathsf T}e_1=0$이고 $R_1$은 위삼각이므로 $(A-\mu I)v_1=(V_mQ_1e_1)(R_1)_{11}$이다. 곧 새 첫 열 $V_mQ_1e_1$은 $(A-\mu I)v_1$의 방향이다. 이동을 차례로 쓰면 $v_1^+$는 $\prod_j(A-\mu_jI)v_1$의 방향이다.

다음으로 Arnoldi 관계에 오른쪽에서 $Q$를 곱하면 $AV^+=V^+H^++f_me_m^{\mathsf T}Q$이다. 각 $Q_j$는 Hessenberg 행렬의 QR에서 나왔으므로 그 자체로 위 Hessenberg이고, $p$개의 곱 $Q$는 아래로 부대각이 $p$개인 띠 행렬이다. 따라서 $e_m^{\mathsf T}Q$의 처음 $k-1$개 성분은 $0$이다. $H^+$는 Hessenberg이므로, 이 관계의 처음 $k$열을 떼어 내면
\[
AV_k^+=V_k^+H_k^++\bigl(h^+_{k+1,k}v_{k+1}^++Q_{mk}f_m\bigr)e_k^{\mathsf T}
\]
이다. 괄호 안의 벡터를 $f_k^+$라 하면, 이것은 $V_k^+$의 열과 직교한다. $v_{k+1}^+$는 정규직교 열 $V^+$의 한 열이고 $f_m$은 $V_m$의 열공간에 직교하기 때문이다. 그러므로 이 관계는 $k$걸음 Arnoldi 관계이고, 23장의 암시적 Q 정리에 의해 같은 시작 벡터에서 만든 Arnoldi 관계와 같다.
\end{proof}

정확한 이동을 쓰면 $H^+_k$의 고유값이 원하는 $k$개의 Ritz 값이 되어 재시작은 원하는 Ritz 벡터들의 공간을 그대로 남기고 나머지를 버린다. 그 뒤로 $p$걸음을 더 늘려 다시 $m$걸음의 관계를 만들고 같은 일을 되풀이한다. 이것이 \term{암시적 재시작 Arnoldi}{implicitly restarted Arnoldi}이고, 대칭이면 암시적 재시작 Lanczos이다\cite{calvetti1994restart}. 실수 비대칭 행렬에서 원하지 않는 Ritz 값이 켤레 복소수 쌍이면 23장의 Francis 이중 이동처럼 두 이동을 한꺼번에 실수 산술로 적용한다.

수학적으로 같은 일을 더 단순하게 하는 방법도 있다. \term{두꺼운 재시작}{thick restart}은 Ritz 벡터들로 직접 재시작한다. Ritz 값을 대각에, 잔차와의 결합을 마지막 행에 둔 \term{화살 꼴}{arrowhead} 사영 행렬에서 다시 확장하는 방법이다\cite{wu2000thick}. 정확한 이동의 암시적 재시작과 같은 공간을 낳는다. 비대칭 행렬에서는 Ritz 벡터 대신 Schur 벡터를 남기는 Stewart의 Krylov--Schur 방법이 같은 역할을 한다\cite{stewart2001krylovschur}. 이미 수렴한 Ritz 쌍은 \term{고정}{locking}하여 더 이상 바꾸지 않고, 남은 부분에서만 재시작하는 \term{수축}{deflation}을 더하면 여러 고유값을 안정적으로 모을 수 있다\cite{lehoucq1996deflation}.

재시작이 얼마나 일을 하는지 보자. $n=1000$의 1차원 Laplace 행렬에서 가장 작은 고유값 $5$개를 상대 잔차 $10^{-10}$까지, 정확한 이동과 같은 두꺼운 재시작 Lanczos로 구하면, 기저의 크기 $m=20$에서는 곱이 $5500$번(재시작 $686$번), $m=40$에서는 $1624$번이 든다. 가장 큰 고유값 $5$개도 $5140$번과 $1642$번으로 비슷하다. 이 행렬의 스펙트럼은 양 끝 모두에서 고유값이 빽빽하기 때문이다. 기저를 두 배로 키우면 곱의 수가 $3.4$배 줄었다. 큰 기저가 좋은 거르개를 만들기 때문이다. 대신 기저의 저장과 직교화 비용이 늘어난다.

\section{스펙트럼 변환}
\label{sec:spectral-transform}

\subsection{이동--역}

\Cref{thm:kps}에 따르면 Krylov 방법은 바깥쪽의 잘 떨어진 고유값을 빨리 찾는다. 그러므로 원하는 고유값을 행렬을 바꾸어 바깥쪽의 잘 떨어진 고유값으로 만들면 된다. \term{이동--역}{shift-invert} 변환은 $A$ 대신
\[
\mathrm{OP}=(A-\sigma I)^{-1}
\]
에 Krylov 방법을 쓴다. $A$의 고유값 $\lambda$는 $\mathrm{OP}$의 고유값 $\nu=1/(\lambda-\sigma)$가 되고 고유벡터는 그대로이다. $\sigma$에 가까운 고유값은 절댓값이 큰 $\nu$가 되고, 그 사이의 간격도 크게 벌어진다. 23장의 역반복과 같은 생각이며, 그 대가는 곱 대신 $(A-\sigma I)x=y$를 푸는 것이다. 보통 26장의 희소 LU로 한 번 분해해 두고 걸음마다 풀이 하나씩을 한다\cite{ericsson1980spectral}.

앞 절의 1차원 Laplace 행렬로 $\sigma=0$의 이동--역을 쓰면, 가장 작은 고유값들은 $\nu=1/\lambda$로 가장 큰 쪽에 간다. 가장 작은 두 고유값의 비가 $\lambda_2/\lambda_1\approx4$이므로 $\nu_2/\nu_1\approx1/4$이 되어, 가장 큰 $\nu$가 나머지에서 멀리 떨어진다. 기저의 크기 $m=20$의 첫 주기, 곧 삼중대각 풀이 $20$번 안에 가장 작은 고유값 $5$개가 상대 $10^{-11}$ 안으로 수렴한다. 곱 $5500$번과 풀이 $20$번의 차이이다. 분해가 가능하다면 안쪽의 고유값은 거의 언제나 이동--역으로 구하는 것이 맞다. 분해가 불가능하면 내부 반복으로 풀거나, Ritz 값 대신 $(A-\sigma I)^{-1}$의 Ritz 값에 해당하는 \term{조화 Ritz 값}{harmonic Ritz value}을 곱만으로 구하는 방법을 쓴다\cite{morgan1991harmonic}.

\subsection{일반화 고유값 문제}

구조 해석이나 진동 문제에서 흔한 것은 대칭 $A$와 대칭 양정치 $M$에 대한 \term{일반화 고유값 문제}{generalized eigenvalue problem} $Ax=\lambda Mx$이다. $M^{-1}A$는 대칭이 아니지만 $M$-내적 $\langle x,y\rangle_M=x^{\mathsf T}My$에 대해 자기수반이다. $\langle M^{-1}Ax,y\rangle_M=x^{\mathsf T}Ay=\langle x,M^{-1}Ay\rangle_M$이기 때문이다. 그래서 Lanczos 과정을 $M$-내적으로 하면 세 항의 점화식과 대칭 삼중대각 행렬이 그대로 살아난다. 이동--역도 $\mathrm{OP}=(A-\sigma M)^{-1}M$으로 같은 일을 한다. 이 밖에 좌굴 문제를 위한 $\mathrm{OP}=(A-\sigma M)^{-1}A$와 Cayley 변환 $\mathrm{OP}=(A-\sigma M)^{-1}(A+\sigma M)$도 쓰는데, 모두 $\lambda$를 $\nu$로 보내는 유리함수가 원하는 고유값을 바깥쪽으로 보내도록 고른 것이다\cite{lehoucq1998arpack}.

\section{특이값}
\label{sec:large-svd}

큰 희소 행렬의 큰 특이값 몇 개, 예컨대 주성분 분석이나 낮은 랭크 근사에 필요한 것은 $A^{\mathsf T}A$의 고유값 문제로 풀고 싶어진다. 그러나 24장에서 본 대로 $A^{\mathsf T}A$를 거치면 작은 특이값이 망가진다. 특이값이 $1$부터 $10^{-12}$까지 로그 간격인 $200\times50$ 행렬을 만들어 견주면, $A^{\mathsf T}A$의 고유값의 제곱근은 $\sigma=2.2\times10^{-5}$에서 상대오차가 $3.4\times10^{-9}$, $\sigma=7.9\times10^{-8}$에서 $1.8\times10^{-3}$이고, $\sigma=10^{-12}$에서는 $9.9\times10^3$이다. 계산한 $A^{\mathsf T}A$는 음의 고유값도 $7$개 가진다. 반면 밀집 SVD는 같은 특이값을 절대오차 $u\sigma_1$ 수준으로 구한다. $A^{\mathsf T}A$의 반올림은 $u\sigma_1^2$ 크기여서 $\sigma<\sqrt u\,\sigma_1$인 특이값을 덮는다.

대안은 앞 장의 LSQR과 같은 Golub--Kahan 이중대각화이다\cite{golub1965kahan}. $A$와 $A^{\mathsf T}$를 번갈아 곱해 $AV_k=U_{k+1}B_k$와 $A^{\mathsf T}U_{k+1}=V_kB_k^{\mathsf T}+\alpha_{k+1}v_{k+1}e_{k+1}^{\mathsf T}$를 만들고, 작은 이중대각 행렬 $B_k$의 특이값 분해 $B_k=U_B\Sigma V_B^{\mathsf T}$에서 Ritz 특이값과 Ritz 특이벡터 $U_{k+1}U_B$, $V_kV_B$를 얻는다. $A^{\mathsf T}A$를 만들지 않으므로 작은 특이값의 정보가 반올림에 덮이지 않고, 24장의 이중대각 행렬의 상대 정확도도 살릴 수 있다. 이 과정은 수학적으로 $A^{\mathsf T}A$에 대한 Lanczos와 같고, Ritz 값이 같은 방식으로 바깥쪽부터 수렴한다. 수학적으로 같은 정보를 갖는 대칭 행렬 $\begin{psmallmatrix}0&A\\A^{\mathsf T}&0\end{psmallmatrix}$의 고유값은 $\pm\sigma_i$이므로 정확도는 지키지만, 작은 특이값은 스펙트럼의 한가운데에 놓여 찾기 어렵다.

Golub--Kahan--Lanczos 과정에도 재시작이 필요하다. Baglama와 Reichel은 두꺼운 재시작을 이중대각화에 맞게 옮겨 원하는 Ritz 특이 삼중쌍을 남기고 $B$를 화살 꼴로 둔 채 다시 확장하는 방법을 만들었다\cite{baglama2005}. 가장 작은 특이값은 이 방법으로도 느리다. 1차원 Laplace 행렬의 가장 작은 고유값이 느렸던 것과 이유가 같다.

\section{stardust의 ARPACK}
\label{sec:arpack-stardust}

큰 고유값 문제의 사실상의 표준은 Lehoucq, Sorensen, Yang의 ARPACK이다\cite{lehoucq1998arpack}. 대칭 문제의 \code{dsaupd}와 비대칭 문제의 \code{dnaupd}가 암시적 재시작 Lanczos와 Arnoldi를 하고, \code{dseupd}와 \code{dneupd}가 수렴한 뒤 Ritz 벡터를 만든다. 사용자는 곱을 직접 하는데, ARPACK이 ``이 벡터에 OP를 곱해 달라''는 요청을 돌려주고 사용자가 결과를 넣어 다시 부르는 \term{역통신}{reverse communication} 방식이다. 그래서 ARPACK은 행렬의 저장 형식을 전혀 모르고, 이동--역의 풀이도 사용자가 한다. scipy의 \code{eigsh}, \code{eigs}는 이 ARPACK을 감싼 것이다.

\begin{instardust}[\code{docs/design/sparse.md} §3.8]
Accelerate에는 ARPACK이 없다. 그래서 stardust는 scipy가 돌리는 ARPACK을 서브루틴 하나에 C++ 함수 하나로 옮겼다. 시작 벡터를 만드는 \code{getv0}, 걸음을 늘리는 \code{saitr}과 \code{naitr}(위의 $\eta=0.717$ 기준의 재직교화를 한 번까지), 이동을 고르는 \code{sgets}와 \code{ngets}, 수렴을 판정하는 \code{sconv}와 \code{nconv}, 암시적 재시작의 QR 걸음인 \code{sapps}와 \code{napps}(대칭은 Givens, 실수 비대칭은 켤레 이동쌍의 Householder 이중 이동), 재시작과 수축을 다루는 \code{saup2}와 \code{naup2}, 그리고 Ritz 벡터를 만드는 \code{seupd}와 \code{neupd}이다. 정렬마저 ARPACK의 셸 정렬을 그대로 옮겼다. 역통신은 연산자의 직접 호출로 바꾸었고, 작은 밀집 문제는 ARPACK이 부르는 LAPACK 함수를 Accelerate에서 부른다. 큰 벡터의 일, 곧 $V^{\mathsf T}x$와 $x-Vh$는 행 $2^{14}$개의 고정 조각마다 BLAS의 행렬--벡터 곱을 부르고 조각의 부분합을 조각 순서로 더하므로, 앞 장의 반복법처럼 스레드 수와 상관없이 같은 비트가 나온다.
\end{instardust}

\begin{instardust}[\code{docs/design/sparse.md} §3.8]
원본 소스를 내려받지 않고 옮겼으므로, 설계 문서는 scipy와 실험으로 견주어 확인한 세부를 적어 둔다. 예컨대 허용오차 \code{tol=0}은 $2^{-52}$를 뜻한다. $2^{-53}$으로 두면 1차원 Laplace 꼴 문제에서 OP의 횟수가 $1441$ 대 $1421$로 scipy와 갈라지고 $2^{-52}$이면 같다. 그런데 재시작에서 분리를 판정하는 기준은 $2^{-53}$에서 나온 것이다. 이렇게 맞춘 뒤 스무 문제에서 열세 문제의 OP 횟수가 scipy와 같았다. 나머지는 시작 벡터의 한 원소를 한 ulp 바꾼 scipy끼리의 퍼짐 안에 있었다. 이 장에서 본 수렴 판정과 재직교화가 반올림에 민감하기 때문이다. 시작 벡터를 주지 않으면 stardust는 고정된 시드의 난수를 써서 스레드 수와 행렬의 저장 형식에 상관없이 같은 입력에 언제나 같은 비트를 낸다.
\end{instardust}

\begin{instardust}[\code{docs/design/sparse.md} §3.8]
\code{svds}는 scipy처럼 $A^{\mathsf T}A$에 \code{eigsh}를 쓰지 않고 Baglama--Reichel의 두꺼운 재시작 Golub--Kahan--Lanczos 이중대각화로 직접 짰다. Ritz 삼중쌍은 $\|A^{\mathsf T}u-\sigma v\|=\|f\|\,|(U_B)_{w,i}|$가 허용오차 안이면 수렴한 것으로 보는데, \eqref{eq:ritz-residual}의 이중대각화 판이다. 설계 문서의 정확도 측정에서 랭크가 부족한 $80\times30$ 행렬의 가장 작은 특이값은 scipy가 계약의 $2\times10^{12}$배나 틀렸고 stardust는 계약의 $0.009$배였다. 위에서 본 $A^{\mathsf T}A$의 반올림이 작은 특이값을 덮어 버렸다.
\end{instardust}

\section{맺음}

큰 고유값 문제는 앞 장과 같은 Krylov 부분공간 위에서 풀린다. 다만 여기서는 해가 하나의 벡터가 아니라 불변 부분공간이어서 기저를 어떻게 지키고 무엇을 남기며 다시 시작하느냐가 핵심이 된다. Lanczos 과정의 직교성 상실은 처음에는 방법의 결함으로 보였지만 Paige의 해명 뒤에는 수렴의 신호로 이해되었다. 암시적 재시작은 23장의 QR 알고리즘을 작은 행렬에 쓰는 것만으로 큰 공간의 시작 벡터를 거르며, 스펙트럼 변환은 다항식으로 닿지 못하는 안쪽 고유값을 유리함수로 바깥에 끌어낸다. 수치선형대수를 다룬 이 부는 여기서 끝나고 다음 부에서는 함수를 다항식과 유리함수로 근사하는 문제로 넘어간다. Chebyshev 다항식은 그곳에서 다시 주인공이 된다.

큰 대칭 고유값 문제는 Parlett의 책\cite{parlett1998}이 고전이고, Lanczos 과정의 유한 정밀도 이론은 Meurant와 Strakoš의 개관\cite{meurant2006}에 정리되어 있다. ARPACK의 사용자 지침서\cite{lehoucq1998arpack}는 암시적 재시작과 스펙트럼 변환의 실제를 자세히 설명한다.

\chapterreferences

